% chapters/ch4-sec4.tex -- 4. Coupling of Dynamic Oscillations to the Trajectory Equations; 4.1-4.4
% (PDF 610-634, printed pp. 578-602), from the heading "4." (mid-PDF 610) up to but not including
% the heading "5." (mid-PDF 634). Equations 132-166b, Tables 1-2.
\section{Coupling of Dynamic Oscillations to the Trajectory Equations}\label{ch4:sec:4}

Having obtained satisfactory solutions to all aspects of point-mass model rocket flight behavior
under the assumption of a uniformly zero angle of attack, we are now in a position to consider the
effects of the rigid-body motion discussed in Chapter~\ref{ch2} in the calculation of altitude
capability.

It would be senseless---and well-nigh impossible---to attempt to present equations for the
detailed calculation of a rocket's altitude performance under conditions of completely general
in-flight disturbances because of the virtually infinite variety of forcing functions which may
cause oscillations of the vehicle. The winds at different altitudes have different speeds and
directions, which moreover change from moment to moment; it is therefore not possible for a
modeler to know precisely the wind profile his vehicle will encounter in any given flight.
Impulsive forcing due to launch departure and staging effects, and the precise nature of the
forcing due to slight propulsive malfunctions such as solid-particle ejection, are also impossible
to predict with any precision. Because of these uncertainties in the nature of the disturbances
encountered in flight, a detailed method for the precise calculation of the maximum altitude
attained by a rocket subjected to completely general forcing functions would be useless. It is a
far more meaningful and fruitful task to solve for the reduction in altitude caused by several
standard and well-defined forcing functions, with the object of \emph{revealing the structural and
aerodynamic characteristics of the vehicle that contribute to the reduction in altitude, thereby
facilitating the design of rockets which are least affected by in-flight disturbances}. It is to
this purpose that the solutions of this section will be directed. Derivations of the solutions
will be carried out for a set of standard forcing functions which form the \emph{basis} for all
possible rigid-body rotative disturbances encountered in model rocket flight, which is to say that
any possible forcing function can be synthesized from some combination of this set of standard
disturbances.

\subsection{The Differential Equations of Motion with Perturbation Terms}\label{ch4:sec:4.1}

In Section~\ref{ch4:sec:1.2} we wrote down the general differential equations of motion which
included perturbation terms due to increases in drag and the reduction of effective thrust due to
pitching and yawing of the vehicle. Later in this section we shall attempt the solution of these
equations. First, however, a review of the assumptions which led to the formation of the vertical
equation of motion~\eqref{ch4:eq:16} is in order. It is important to remember that
equation~\eqref{ch4:eq:16} is \emph{not} the exact differential equation of vertical motion; it
is an approximation. The dynamic analysis of Chapter~\ref{ch2} is valid \emph{only} provided that
the model's \CG{} does not move laterally too much during any given disturbance. In actuality, the
momentary deflections of an oscillating rocket will cause it to head in a direction that is not
tangent to the intended direction of flight and the \CG{} \emph{will} move slightly to the side
before the angle of attack passes through zero on its way toward the opposite extreme of the
oscillation. \emph{This effect will contribute to the reduction in altitude} caused by the
disturbance and will appear as an additional term not originally included in
equation~\eqref{ch4:eq:16}. The treatments of this section will, however, be restricted to cases
for which the dynamic analyses of Chapter~\ref{ch2} are valid. It will therefore be assumed that
the model's frequency of oscillation is sufficiently high, and its amplitude sufficiently small,
that the oscillatory horizontal translations of the \CG{} are negligibly slight. Such an
assumption is justified, since (as was shown in Chapter~\ref{ch2}) a too-low oscillation frequency
or a critically-damped or overdamped response is incompatible with dynamically stable flight.

It will thus be assumed that the rocket's \CG{} is travelling in a purely vertical direction and
develops no lateral components of velocity. Under such an assumption,
equation~\eqref{ch4:eq:16} becomes an exact description of the rocket's vertical motion. We are
then free to consider the method by which we will solve~\eqref{ch4:eq:16}.

Since a stable rocket will, as shown in Chapter~\ref{ch2}, practically never experience pitch or
yaw angles exceeding $12\dg$, it is possible to expand some of the terms of
equation~\eqref{ch4:eq:16} containing functions of $\alpha(t)$ as factors, in powers of
$\alpha(t)$, where $\alpha(t)$ is the generalized yaw-pitch angle which equals
\begin{equation}\label{ch4:eq:132}
\alpha(t) = \sqrt{[\alpha_x(t)]^{2} + [\alpha_y(t)]^{2}}
\end{equation}
and $\alpha_x(t)$ and $\alpha_y(t)$ are the yaw and pitch angles of the model as functions of
time.

As stated in Chapter~\ref{ch1} and confirmed in greater detail in Chapter~\ref{ch3}, the function
$f(\alpha)$ appearing in~\eqref{ch4:eq:16} can be represented to a high order of accuracy by a
quadratic function in $\alpha(t)$, so that the approximate drag as a function of the angle of
attack can be written as
\begin{equation}\label{ch4:eq:133}
\text{Drag} = (k + \epsilon\alpha^{2})v^{2}
\end{equation}
where the determination of the constants $k$ and $\epsilon$ has been discussed in
Chapters~\ref{ch1} and~\ref{ch3}, and $\alpha$ is just $\alpha(t)$, the generalized angle of
attack. It is also possible to replace $\cos(\alpha)$ by its power series expansion
\begin{equation}\label{ch4:eq:134}
\cos(\alpha) = 1 - \frac{\alpha^{2}}{2} + \frac{\alpha^{4}}{4} - \ldots
\end{equation}
and to truncate the expansion after its $\alpha^{2}$ term, since $\alpha$ will never exceed 0.2
radian for cases amenable to analysis by the methods of Chapter~\ref{ch2}:
\begin{equation}\label{ch4:eq:135}
\cos(\alpha) \cong 1 - \frac{\alpha^{2}}{2}
\end{equation}

Equation~\eqref{ch4:eq:16} may then be written as
\begin{equation}\label{ch4:eq:136}
m(t)\frac{dv}{dt} = F(t)\left[1 - \frac{\alpha^{2}}{2}\right] - m(t)g - (k + \epsilon\alpha^{2})v^{2}
\end{equation}
where the $y$ subscript has been discarded as unnecessary due to the fact that we are here
considering cases of vertical flight only.

Despite the simplifying approximations made to obtain it, equation~\eqref{ch4:eq:136} is
sufficiently complicated that no practical, closed-form solution to it has yet been found. As in
Section~\ref{ch4:sec:3}, therefore, we shall have to resort to numerical solutions obtained by
incremental methods similar to those of Section~\ref{ch4:sec:2.4}. While it is not possible to
visualize the sensitive parameters influencing a model's behavior as concisely when such
computer-oriented techniques are used as it would be if an accurate and reasonably simple
closed-form solution were available, it will be found that one \emph{can} gain by example some
feeling for the effects of various disturbances on the altitude capability of representative model
rockets. Such numerical data, interpreted in the light of the analytical results of
Chapter~\ref{ch2}, enable us to fulfill our stated purpose of extracting criteria of use in
designing models which possess maximum resistance to in-flight perturbations.

\subsection{A Numerical Method for the Digital Computation of Altitude Performance in Cases of
Oscillating Rockets}\label{ch4:sec:4.2}

Like the vertical equation of motion for non-oscillating rockets, equation~\eqref{ch4:eq:136} can
be solved to any desired accuracy using numerical techniques in which the differential time $dt$
is replaced by a finite interval $\Delta t$. In order to obtain solutions for rockets undergoing
rigid-body oscillations, however, it is necessary to add some additional computing steps derived
from Euler's dynamical equations---equations~\eqref{ch2:eq:14} of Chapter~\ref{ch2}. Listed below
is the computational procedure for completing a single iterative loop analogous to the
drag-from-prior-velocity method of Section~\ref{ch4:sec:2.4}:
% (137)-(153) are one printed run (PDF 614-615, STYLE.md section 15): the statements (137)-(143)
% start at one column, the "=" signs of (144)-(153) are lined up.
\begin{align}
&\Delta v = \Delta t\left\{F(t)\left[1 - \frac{\alpha^{2}}{2}\right] - m(t)g
  - (k + \epsilon\alpha^{2})v^{2}\right\}\Big/m(t) \label{ch4:eq:137}\\
&\Delta y = \Delta t\left(v + \frac{\Delta v}{2}\right) \label{ch4:eq:138}\\
&v = v + \Delta v \label{ch4:eq:139}\\
&y = y + \Delta y \label{ch4:eq:140}\\
&\text{calculate } C_1 \label{ch4:eq:141}\\
&\text{calculate } C_2 \label{ch4:eq:142}\\
&\text{calculate or input } \omega_z \label{ch4:eq:143}\\
\Delta\omega_x &= \Delta t[f_x(t) - C_2\omega_x - C_1\alpha_x - I_R\omega_y\omega_z]/I_L
  \label{ch4:eq:144}\\
\Delta\omega_y &= \Delta t[f_y(t) - C_2\omega_y - C_1\alpha_y + I_R\omega_y\omega_z]/I_L
  \label{ch4:eq:145}\\
\Delta\alpha_x &= \Delta t\left[\omega_x + \frac{\Delta\omega_x}{2}\right] \label{ch4:eq:146}\\
\Delta\alpha_y &= \Delta t\left[\omega_y + \frac{\Delta\omega_y}{2}\right] \label{ch4:eq:147}\\
\omega_x &= \omega_x + \Delta\omega_x \label{ch4:eq:148}\\
\omega_y &= \omega_y + \Delta\omega_y \label{ch4:eq:149}\\
\alpha_x &= \alpha_x + \Delta\alpha_x \label{ch4:eq:150}\\
\alpha_y &= \alpha_y + \Delta\alpha_y \label{ch4:eq:151}\\
\alpha &= \sqrt{\alpha_x^{2} + \alpha_y^{2}} \label{ch4:eq:152}\\
t &= t + \Delta t \label{ch4:eq:153}
\end{align}
Again, as in Section~\ref{ch4:sec:2.4}, expressions such as (175) are not equations as we
understand them, but arithmetic assignment statements which instruct the computer to replace the
``old'' or ``former'' value of a variable with its ``new'' or ``present'' value. The notation of
Chapter~\ref{ch2} has been used for all the expressions involving rigid-body dynamics: subscripts
$x$ and $y$ \emph{do not} refer to the Cartesian axes used in the general equations of nonvertical
motion, but rather to the pitch and yaw axes of the vehicle; $\omega$ denotes angular velocity;
$\alpha$ denotes angular displacement or angle of attack; the $f(t)$ are forcing functions or
perturbing moments that may vary with time; $C_1$ is the corrective moment coefficient; $C_2$ is
the damping moment coefficient; $I_L$ is the longitudinal moment of inertia; and $I_R$ is the
radial moment of inertia.

You should note that equations~\eqref{ch4:eq:144} and~\eqref{ch4:eq:145} are approximate in that
they assume constant, average values for $I_L$ and $I_R$ (the moments of inertia actually change
slightly as the propellant is expended). Exact, time-varying values could be used---as could the
exact expressions for thrust reduction and drag increase due to angle of attack---but the increase
in accuracy obtained thereby is not worth the consequent increase in complexity. The computer is
going to have to perform the calculations~\eqref{ch4:eq:137} through~\eqref{ch4:eq:153} a great
many times, so even a slight time saving in a few steps of the loop can add up to worthwhile
savings in time and expense over the full course of a computer ``run''.

Steps~\eqref{ch4:eq:141}, \eqref{ch4:eq:142}, and~\eqref{ch4:eq:143} require the use of formulae
developed in Chapter~\ref{ch2}. Specifically, $C_1$ is proportional to the square of the velocity
according to the relation expressed in equation~\eqref{ch2:eq:95} of Chapter~\ref{ch2}, while
$C_2$ depends on both the velocity and the motor characteristics as given in
equations~\eqref{ch2:eq:97} through~\eqref{ch2:eq:101} of that chapter. In cases where the roll
rate $\omega_z$ is induced by canted fins it may be calculated from equations~\eqref{ch2:eq:115}
through~\eqref{ch2:eq:117} of Chapter~\ref{ch2}; if induced by other means it must be determined
by a formula appropriate to the method of roll forcing. If the assumption of zero roll rate is to
be made the value $\omega_z = 0$ must be input (or you can write a special program for use with
non-rolling rockets only, in which the equations assume $\omega_z$ equal to zero).

As you can determine by working a few numerical examples, the angular velocities---called pitch,
yaw, and roll \emph{rates} by engineers---can become quite high for model rockets near burnout
velocity. The method of theoretical velocity increments presented in Section~\ref{ch4:sec:2.4}
therefore loses its computational advantage over the drag-from-prior-velocity method when applied
to cases of oscillating rockets. That advantage consisted of a tolerance for relatively large time
increments while maintaining good accuracy, and it is lost when oscillating rockets are considered
because short intervals are required for an accurate determination of the angle of attack. The
modification of the theoretical-velocity-increment method for oscillating vehicles will not,
therefore, be presented here. For best results, $\Delta t$ as used in~\eqref{ch4:eq:137}
through~\eqref{ch4:eq:153} should be of the order of 0.001 second---and it is unlikely that the
theoretical-velocity-increment method would tolerate much more than this. For hand calculations it
is of course impractical to attempt the use of $\Delta t$'s less than 0.1 second. Those who perform
such calculations should not, therefore, expect too much in the way of accuracy from the procedure
detailed above.

\subsection{Solutions for Standard Forcing Functions}\label{ch4:sec:4.3}

In order to apply the method of Section~\ref{ch4:sec:4.2} in computing the altitude performance of
a model rocket perturbed by one of the standard forcing functions described in Chapter~\ref{ch2} it
is necessary to select specific functional forms for $f_x(t)$ and $f_y(t)$ and specific initial
values for $\omega_x$, $\omega_y$, $\alpha_x$, and $\alpha_y$. The remainder of this section will
present such formulae and values for use in determining the effects of a homogeneous response and
of responses to each of the standard forcing functions considered in Chapter~\ref{ch2}. Our
treatment will be restricted to single-staged models (although the method of
Section~\ref{ch4:sec:4.2} is completely general and can be applied to rockets of any number of
stages), since all information of design interest to model rocketeers concerning dynamic response
is obtainable from the study of single-stage cases.

\subsubsection{Homogeneous Response for General Initial Conditions}\label{ch4:sec:4.3.1}

When one is primarily concerned with the behavior of a model once it has been placed in a given
state of angular displacement and angular velocity, without regard to the precise nature of the
disturbance that caused it to assume that state, it is best to compute the \emph{homogeneous
response} resulting from the initial conditions that have been assumed. For such a calculation we
assume that, prior to some time $t_o$ after liftoff, the vehicle has been flying vertically
straight and true. For values of $t$ less than $t_o$, therefore,
% The third member of this group is printed "(156c)" (evidently 154c; the plain (156) follows on
% PDF 619): the Chapter 3 (203)-(204) pattern, with \refstepcounter giving the group label (154)
% and advancing the counter so that the next group is (155).
\refstepcounter{equation}\label{ch4:eq:154}%
\begin{empheq}[right={\empheqrbrace\quad t < t_o}]{align*}
\alpha_x &= \alpha_y = 0 \tag{154a}\label{ch4:eq:154a}\\
\omega_x &= \omega_y = 0 \tag{154b}\label{ch4:eq:154b}\\
f_x(t) &= f_y(t) = 0 \tag{156c}\label{ch4:eq:156c}
\end{empheq}
Equations~\eqref{ch4:eq:154} are said to specify the rocket's state prior to $t_o$ as
\emph{rotationally quiescent}; it follows that the iterative procedure described in
Section~\ref{ch4:sec:4.2} produces values of $v$ and $y$ identical to those obtained by the methods
of Section~\ref{ch4:sec:2.4} for all values of $t$ up through $t_o$.

In the instant before $t_o$, it is assumed that some unspecified occurrence disturbs the rocket
from a straight-and-true attitude, leaving it with nonzero angular displacements and velocities as
follows:
\begin{subequations}\label{ch4:eq:155}
\begin{empheq}[right={\empheqrbrace\quad t = t_o}]{align}
\alpha_x &= \alpha_{xo} \label{ch4:eq:155a}\\
\alpha_y &= \alpha_{yo} \label{ch4:eq:155b}\\
\omega_x &= \omega_{xo} \label{ch4:eq:155c}\\
\omega_y &= \omega_{yo} \label{ch4:eq:155d}
\end{empheq}
\end{subequations}
where the zero-subscripted quantities are the initial values of angular displacement and velocity
to be considered. These values must be chosen by the person performing the calculation and are
arbitrary except for the restriction that they cannot be so large as to cause the model to have an
angle of attack greater than 0.2 radian at any time during the subsequent oscillation. Since a
homogeneous response is defined as a response for which there is no forcing function, the yaw and
pitch forcing remain uniformly zero for all values of time, after as well as before $t_o$:
\begin{equation}\label{ch4:eq:156}
f_x(t) = f_y(t) = 0 \qquad (t \ge t_o)
\end{equation}
Computations of this kind may be performed for either zero or nonzero roll rate. In either case, no
generality will be lost and the response behavior will be more easily interpreted if you assume
nonzero initial angular displacement and angular velocity about one axis only; it is merely for the
sake of completeness that I have taken into account both the pitch and the yaw axis in setting
initial conditions.

\subsubsection{Step Response}\label{ch4:sec:4.3.2}

In order to compute the response of a model to step forcing that suddenly arises at some time $t_o$
into the flight we again assume that the rocket is in a rotationally quiescent state for all time
prior to $t_o$; equations~\eqref{ch4:eq:154} thus remain valid for this problem. At the moment the
step forcing begins the model will still be in an undisturbed state, so the initial conditions are
\begin{subequations}\label{ch4:eq:157}
\begin{empheq}[right={\empheqrbrace\quad t = t_o}]{align}
\alpha_x &= \alpha_y = 0 \label{ch4:eq:157a}\\
\omega_x &= \omega_y = 0 \label{ch4:eq:157b}
\end{empheq}
\end{subequations}
The forcing functions, however, become nonzero at this instant and persist as long as the step is
assumed to last---i.e., until some time $t_1$ at which the forcing ``steps down'' again to zero (or
until flight apex occurs, which is equivalent to setting $t_1$ equal to the duration of the upward
flight). Then
\begin{subequations}\label{ch4:eq:158}
\begin{empheq}[right={\empheqrbrace\quad t_o \le t \le t_1}]{align}
f_x(t) &= M_x \label{ch4:eq:158a}\\
f_y(t) &= M_y \label{ch4:eq:158b}
\end{empheq}
\end{subequations}
where $M_x$ and $M_y$ are the disturbing \emph{torques}, or moments, about the yaw and pitch axes,
respectively, and
\begin{equation}\label{ch4:eq:159}
f_x(t) = f_y(t) = 0 \qquad (t \ge t_1)
\end{equation}
In cases of non-spinning rockets where the step forcing is due to slight aerodynamic asymmetries,
the strength of the perturbing moments will vary as the square of the airspeed. Where the step
arises from a uniform horizontal wind the magnitudes of $M_x$ and $M_y$ will be constant throughout
the time during which the wind is blowing.

It is important to note that step forcing is \emph{not} an accurate representation of the action
of a horizontal wind upon the model except for relatively short times---on the order of a few
tenths of a second for most model rockets. You may remember from the discussion of step responses
in Chapter~\ref{ch2} that true step forcing applied for a sufficiently long time causes the vehicle
to assume an angle of attack equal to $M/C_1$, where $M$ is the disturbing moment. A horizontal
wind, however, adds vectorially to the model's airspeed and \emph{changes the direction} in which
the rocket must fly to be at \emph{zero} angle of attack; its long-term effect is therefore to make
the model tip over into a nonvertical trajectory (we as model rocketeers are used to referring to
this as \emph{weathercocking}), lowering its altitude by this means as well as by the decreased
effective thrust and increased drag due to the oscillatory response. The steady angle of attack
predicted by the step-response treatment produces a considerable overestimate of the aerodynamic
drag on the perturbed rocket, though, and this drag overestimate nearly always more than
compensates for the failure to consider the nonverticality of the flight path. For this reason it
is acceptable to use step forcing to simulate horizontal wind effects in making performance
degradation estimates despite the admitted analytical inaccuracy of the technique.

In making step response calculations you may choose $M_x$ and $M_y$ arbitrarily, subject only to
the restriction that they be small enough to make sure the rocket will not exceed an angle of
attack of 0.2 radian at any time during the response. It admittedly takes some numerical
experimentation to establish such upper bounds, but the work is unavoidable if the validity of
Chapter~\ref{ch2}'s dynamic analyses is to be guaranteed. By reasoning similar to that of
Section~\ref{ch4:sec:4.3.1} it may be concluded that no information regarding the response will be
lost if one of the disturbing torques is considered to be zero; this is true whether or not the
rocket is spinning.

\subsubsection{Impulse Response}\label{ch4:sec:4.3.3}

In order to determine a model's response to impulsive forcing occurring at the instant $t_o$ after
liftoff we again assume the quiescent prior rotational state specified by equations (189). The
effect of the impulse, as discussed in Chapter~\ref{ch2}, will then be to produce a homogeneous
response for time greater than $t_o$ with the following specialized set of initial conditions:
\begin{subequations}\label{ch4:eq:160}
\begin{empheq}[right={\empheqrbrace\quad t = t_o}]{align}
\alpha_x &= \alpha_y = 0 \label{ch4:eq:160a}\\
\omega_x &= \frac{H_x}{I_L} \label{ch4:eq:160b}\\
\omega_y &= \frac{H_y}{I_L} \label{ch4:eq:160c}
\end{empheq}
\end{subequations}
where $H_x$ and $H_y$ are the components of the impulse about the yaw and pitch axes,
respectively.

Because of the unique nature of impulsive forcing (the impulse is a mathematical idealization
involving the assumption of an infinite disturbing moment applied for an infinitesimal time
period), an impulse response must be computed by setting the yaw and pitch forcing functions to
zero for all time:
\begin{equation}\label{ch4:eq:161}
f_x(t) = f_y(t) = 0
\end{equation}
Actually, $f_x(t)$ and $f_y(t)$ are both effectively infinite at the time $t_o$, but the results of
this have already been taken into account in equations (195).

Again, any values for $H_x$ and $H_y$---including zero for either one---will be acceptable as long
as these values do not cause the overall angle of attack to exceed 0.2 radian during the subsequent
response.

\subsubsection{Response to Sinusoidal Forcing}\label{ch4:sec:4.3.4}

Although a wide variety of sinusoidal disturbances could be hypothesized for cases of both rolling
and non-rolling rockets, the model rocketeer will find that all the information relating to
sinusoidal forcing that is of interest in designing high-performance models can be obtained by
analyzing two specialized cases:
\begin{enumerate}[label=(\alph*)]
  \item Sinusoidal forcing of a non-rolling rocket about one axis only, in which the amplitude of
    the forcing varies as the square of the airspeed and its frequency varies directly with the
    airspeed; and
  \item Sinusoidal forcing of a rolling rocket in which the sinusoidal nature of the forcing is due
    to the roll rate itself, so that sinusoidal moments exist about the pitch and yaw axes of
    angular frequency equal to the roll rate and of amplitude proportional to the square of the
    airspeed.
\end{enumerate}

A disturbance of the first type may be characterized by the following set of forcing functions:
\begin{subequations}\label{ch4:eq:162}
\begin{align}
f_x(t) &= A_f\sin\omega_f t \label{ch4:eq:162a}\\
f_y(t) &= 0 \label{ch4:eq:162b}
\end{align}
\end{subequations}
where the disturbance has been assumed to occur about the yaw axis. The proportionality of the
forcing amplitude $A_f$ and the forcing frequency $\omega_f$ to the rocket's velocity can be
expressed as
\begin{subequations}\label{ch4:eq:163}
\begin{equation}\label{ch4:eq:163a}
A_f = A_o v^{2}
\end{equation}
and
\begin{equation}\label{ch4:eq:163b}
\omega_f = \omega_o v
\end{equation}
\end{subequations}
where $A_o$ and $\omega_o$ are constants of proportionality which may be thought of as the values
of $A_f$ and $\omega_f$ at an airspeed of one meter per second.

A disturbance of the second type exhibits forcing functions of the form
\begin{subequations}\label{ch4:eq:164}
\begin{align}
f_x(t) &= A_f\sin\omega_z t \label{ch4:eq:164a}\\
f_y(t) &= A_f\cos\omega_z t \label{ch4:eq:164b}
\end{align}
\end{subequations}
Since spin in model rockets is almost invariably induced by aerodynamic means, the roll rate may
be considered linearly proportional to velocity for most cases of interest. The same sort of
proportionality scheme used for the non-rolling case can then be used here:
\begin{subequations}\label{ch4:eq:165}
\begin{align}
A_f &= A_o v^{2} \label{ch4:eq:165a}\\
\omega_z &= \omega_{zo} v \label{ch4:eq:165b}
\end{align}
\end{subequations}
Equation (200a)\ednote{No equation carries the number (200a); equation~\eqref{ch4:eq:165a} is
evidently meant.} applies whether the disturbance is due to aerodynamic or to inertial causes,
since the strength of the aerodynamic perturbation is proportional to the square of the airspeed
and that of the inertial perturbation is proportional to the square of the roll rate, which in turn
varies directly with airspeed.

Both types of perturbations are assumed to arise immediately upon liftoff and to persist
throughout the flight, as such a mathematical model most accurately represents the actual behavior
of sinusoidal forcing in the vast majority of cases of interest. Calculations should therefore be
started at $t = 0$ with rotationally quiescent initial conditions:
\begin{subequations}\label{ch4:eq:166}
\begin{empheq}[right={\empheqrbrace\quad t = 0}]{align}
\alpha_x &= \alpha_y = 0 \label{ch4:eq:166a}\\
\omega_x &= \omega_y = 0 \label{ch4:eq:166b}
\end{empheq}
\end{subequations}

It should be noted that the calculations presented here differ somewhat from the analyses of
Chapter~\ref{ch2}'s Sections~\ref{ch2:sec:3.1.4} and~\ref{ch2:sec:3.2.4} in that we are considering
the \emph{complete} response to sinusoidal forcing rather than the \emph{steady-state} response
only; our numerical calculations will therefore pick up so-called \emph{starting transients} which
are not considered in the discussion of Chapter~\ref{ch2}. The basic character of the response will
not, however, be greatly altered and it is very nearly correct to consider the method of this
section the numerical equivalent to the closed-form analyses of Chapter~\ref{ch2}.

\subsection{The Effect of Dynamic Oscillations on the Altitude Performance of a Typical Model
Rocket}\label{ch4:sec:4.4}

In order to give the reader some ``feel'' for the magnitude of the performance degradations
resulting from the various in-flight disturbances that cause pitching and yawing oscillations in
model rockets, this section will present the results of sample calculations based on
representative disturbances encountered by a typical model rocket. The trial rocket illustrated in
Figure~\ref{ch3:fig:37} of Chapter~\ref{ch3} was chosen as a representative design; reasonable
values of inertial and propulsive characteristics were then supplied which, in addition to the
drag parameters and normal force coefficients computable from the model's geometry, permitted a
complete characterization of the vehicle for computational purposes. A complete list of the
vehicle parameters used in solving the coupled equations of motion by the incremental method of
Section~\ref{ch4:sec:4.2} appears in Table~\ref{ch4:tab:1}.

\begin{table}[htbp]
  \centering
  \caption{Ballistic and dynamic characteristics of the rocket shown in Figure~\ref{ch3:fig:37}
    of Chapter~\ref{ch3}.}
  \label{ch4:tab:1}
  \begin{tabular}{@{}l@{\qquad}l@{}}
    \toprule
    Characteristic & Value \\
    \midrule
    Liftoff mass & 40 g = 0.04 kg \\
    Burnout mass & 31.67 g = 0.03167 kg \\
    Engine thrust & 4.0 N (average) \\
    Burning time & 1.20 sec \\
    Propellant mass & 8.33 g = .00833 kg \\
    Static stability margin & 2.06 cm = 1.0 caliber \\
    $C_1$ & 396$v^{2}$ dyn-cm for $v$ in m/sec \\
    $C_{2A}$ & 242$v$ dyn-cm-sec for $v$ in m/sec \\
    $C_{2R}$ & 661 dyn-cm-sec (average) \\
    $I_L$ & 3640 g-cm$^{2}$ (average) \\
    $I_R$ & 42.4 g-cm$^{2}$ (average) \\
    $\zeta$ & 0.101 neglecting jet damping \\
    $\zeta_c$ & 0.1004 neglecting jet damping \\
    $\omega_{\mathrm{res}}$ & 0.327$v$ rad/sec for $v$ in m/sec \\
    $\omega_{\mathrm{cres}}$ & 0.325$v$ rad/sec for $v$ in m/sec \\
    $(\CDo)_{FB}$ & 0.45 without fin cant \\
    \begin{tabular}[c]{@{}l@{}}Fin cant angle for\\ $\omega_z = \omega_{\mathrm{cres}}$\end{tabular}
      & 0.01945 rad = $1.11\dg$ \\
    \begin{tabular}[c]{@{}l@{}}Fin cant angle for\\ $\omega_z = 10\omega_{\mathrm{cres}}$\end{tabular}
      & 0.1945 rad = $11.1\dg$ \\
    $k$ & \begin{tabular}[t]{@{}l@{}}$0.92 \times 10^{-4}$ kg/m without fin cant\\
      $0.924 \times 10^{-4}$ kg/m with $1.11\dg$ cant\\
      $1.288 \times 10^{-4}$ kg/m with $11.1\dg$ cant\end{tabular} \\
    $\epsilon$ & $37.8 \times 10^{-4}$ (kg/m)/rad$^{2}$ \\
    \bottomrule
  \end{tabular}
\end{table}

Table~\ref{ch4:tab:2} summarizes the results of calculations performed for various cases of input
disturbances. Although these results admittedly apply only to the particular rocket described in
Table~\ref{ch4:tab:1}, they do---in combination with the discussion of Chapter~\ref{ch2}---permit
us to draw certain conclusions of general interest and applicability to model rocket design.

% Table 2 is printed in two row-matched halves (PDF 628-629: columns 1-4; PDF 630-631: columns
% 5-10; caption on PDF 629); it is set as one 10-column longtable (STYLE.md section 15).
\begingroup\small
\begin{longtable}{@{}>{\raggedright\arraybackslash}p{4.1cm}cc
    S[table-format=1.3]S[table-format=3.1]S[table-format=2.1]S[table-format=1.2]
    S[table-format=3.1]S[table-format=3.1]S[table-format=2.2]@{}}
  \caption{The effect of various disturbances on the performance of the rocket shown in
    Figure~\ref{ch3:fig:37} of Chapter~\ref{ch3}. In all cases where disturbances or roll rates are
    velocity-dependent the velocity is given in meters per second, and the amplitudes of
    sinusoidal forcing functions are given in dyn-cm. $t_o$ is the time at which the disturbance
    begins; $\alpha_{\max}$ is the maximum angle of attack produced by the disturbance.}
  \label{ch4:tab:2}\\
  \toprule
  \begin{tabular}[b]{@{}l@{}}Description of disturbance\\ or response\end{tabular}
    & \begin{tabular}[b]{@{}c@{}}$\omega_z$\\ (rad/sec)\end{tabular}
    & \begin{tabular}[b]{@{}c@{}}$t_o$\\ (sec)\end{tabular}
    & {\begin{tabular}[b]{@{}c@{}}$\alpha_{\max}$\\ (rad)\end{tabular}}
    & {\begin{tabular}[b]{@{}c@{}}$v_b$\\ (m/sec)\end{tabular}}
    & {\begin{tabular}[b]{@{}c@{}}$y_b$\\ (m)\end{tabular}}
    & {\begin{tabular}[b]{@{}c@{}}$t_c$\\ (sec)\end{tabular}}
    & {\begin{tabular}[b]{@{}c@{}}$y_c$\\ (m)\end{tabular}}
    & {\begin{tabular}[b]{@{}c@{}}$y_{\max}$\\ (m)\end{tabular}}
    & {\begin{tabular}[b]{@{}c@{}}Percent\\ reduction\\ in $y_{\max}$\end{tabular}} \\
  \midrule
  \endfirsthead
  \multicolumn{10}{@{}l}{\tablename~\thetable\ (continued)}\\
  \toprule
  \begin{tabular}[b]{@{}l@{}}Description of disturbance\\ or response\end{tabular}
    & \begin{tabular}[b]{@{}c@{}}$\omega_z$\\ (rad/sec)\end{tabular}
    & \begin{tabular}[b]{@{}c@{}}$t_o$\\ (sec)\end{tabular}
    & {\begin{tabular}[b]{@{}c@{}}$\alpha_{\max}$\\ (rad)\end{tabular}}
    & {\begin{tabular}[b]{@{}c@{}}$v_b$\\ (m/sec)\end{tabular}}
    & {\begin{tabular}[b]{@{}c@{}}$y_b$\\ (m)\end{tabular}}
    & {\begin{tabular}[b]{@{}c@{}}$t_c$\\ (sec)\end{tabular}}
    & {\begin{tabular}[b]{@{}c@{}}$y_c$\\ (m)\end{tabular}}
    & {\begin{tabular}[b]{@{}c@{}}$y_{\max}$\\ (m)\end{tabular}}
    & {\begin{tabular}[b]{@{}c@{}}Percent\\ reduction\\ in $y_{\max}$\end{tabular}} \\
  \midrule
  \endhead
  \bottomrule
  \endlastfoot
  No disturbance
    & 0 & --- & 0.000 & 111.4 & 78.1 & 6.46 & 265.5 & 343.6 & 0.00 \\[2pt]
  Homogeneous response to\newline $\alpha_{xo} = 0.2$ rad
    & 0 & 0.1 & 0.200 & 110.3 & 76.9 & 6.44 & 262.8 & 339.6 & 1.17 \\[2pt]
  Homogeneous response to\newline $\alpha_{xo} = 0.2$ rad
    & 0 & 1.2 & 0.200 & 111.4 & 78.1 & 6.35 & 254.3 & 332.4 & 3.26 \\[2pt]
  Step of intensity\newline 39.6$v^{2}$ dyn-cm
    & 0 & 0.1 & 0.155 & 105.3 & 75.7 & 5.67 & 211.0 & 286.7 & 16.58 \\[2pt]
  Step of intensity\newline 51.0$v^{2}$ dyn-cm
    & 0 & 0.1 & 0.200 & 101.7 & 74.2 & 5.28 & 186.1 & 260.3 & 24.25 \\[2pt]
  Step of intensity\newline 39.6$v^{2}$ dyn-cm
    & 0 & 1.2 & 0.177 & 111.4 & 78.1 & 5.75 & 220.4 & 298.4 & 13.16 \\[2pt]
  Step of intensity\newline 45.0$v^{2}$ dyn-cm
    & 0 & 1.2 & 0.200 & 111.4 & 78.1 & 5.59 & 210.3 & 288.4 & 16.08 \\[2pt]
  Impulse of strength\newline 2510 dyn-cm-sec
    & 0 & 0.1 & 0.062 & 111.3 & 77.9 & 6.46 & 265.2 & 343.1 & 0.15 \\[2pt]
  Impulse of strength\newline 8080 dyn-cm-sec
    & 0 & 0.1 & 0.200 & 109.9 & 76.7 & 6.43 & 261.9 & 338.6 & 1.46 \\[2pt]
  Impulse of strength\newline 26,600 dyn-cm-sec
    & 0 & 1.2 & 0.175 & 111.4 & 78.1 & 6.35 & 254.7 & 332.8 & 3.14 \\[2pt]
  Impulse of strength\newline 30,400 dyn-cm-sec
    & 0 & 1.2 & 0.200 & 111.4 & 78.1 & 6.32 & 251.5 & 329.5 & 4.10 \\[2pt]
  Sinusoidal forcing of\newline $A_f = 15.9v^{2}$, $\omega_f = \omega_{\mathrm{res}}$
    & 0 & 0.1 & 0.143 & 111.1 & 77.9 & 6.33 & 257.8 & 335.7 & 2.30 \\[2pt]
  Sinusoidal forcing of\newline $A_f = 26.7v^{2}$, $\omega_f = \omega_{\mathrm{res}}$
    & 0 & 0.1 & 0.200 & 110.5 & 77.6 & 6.08 & 242.9 & 320.5 & 6.72 \\[2pt]
  Homogeneous response to\newline $\alpha_{xo} = 0.2$ rad
    & 3.25$v$ & 0.1 & 0.200 & 105.0 & 74.9 & 5.68 & 211.2 & 286.1 & 16.74 \\[2pt]
  Homogeneous response to\newline $\alpha_{xo} = 0.2$ rad
    & 3.25$v$ & 1.2 & 0.200 & 106.0 & 76.0 & 5.60 & 204.6 & 280.6 & 18.34 \\[2pt]
  Step of intensity\newline 39.6$v^{2}$ dyn-cm
    & 3.25$v$ & 0.1 & 0.155 & 100.6 & 73.8 & 5.13 & 176.9 & 250.7 & 27.05 \\[2pt]
  Step of intensity\newline 51.0$v^{2}$ dyn-cm
    & 3.25$v$ & 0.1 & 0.200 & 97.4 & 72.5 & 4.84 & 159.1 & 231.6 & 32.60 \\[2pt]
  Step of intensity\newline 39.6$v^{2}$ dyn-cm
    & 3.25$v$ & 1.2 & 0.176 & 106.0 & 76.1 & 5.20 & 183.8 & 259.8 & 24.39 \\[2pt]
  Step of intensity\newline 45.0$v^{2}$ dyn-cm
    & 3.25$v$ & 1.2 & 0.200 & 106.0 & 76.1 & 5.08 & 176.8 & 252.9 & 26.41 \\[2pt]
  Impulse of strength\newline 2510 dyn-cm-sec
    & 3.25$v$ & 0.1 & 0.062 & 105.9 & 75.9 & 5.69 & 212.9 & 288.8 & 15.95 \\[2pt]
  Impulse of strength\newline 8080 dyn-cm-sec
    & 3.25$v$ & 0.1 & 0.200 & 104.7 & 74.7 & 5.67 & 210.6 & 285.4 & 16.95 \\[2pt]
  Impulse of strength\newline 26,600 dyn-cm-sec
    & 3.25$v$ & 1.2 & 0.175 & 106.0 & 76.1 & 5.61 & 205.0 & 281.0 & 18.23 \\[2pt]
  Impulse of strength\newline 30,400 dyn-cm-sec
    & 3.25$v$ & 1.2 & 0.200 & 106.0 & 76.1 & 5.58 & 202.5 & 278.6 & 18.92 \\[2pt]
  Coupled sinusoidal forcing\newline of $A_f = 15.9v^{2}$, $\omega_z = \omega_{\mathrm{cres}}$
    & 0.325$v$ & 0.1 & 0.141 & 110.7 & 77.7 & 6.19 & 249.5 & 327.2 & 4.78 \\[2pt]
  Coupled sinusoidal forcing\newline of $A_f = 26.7v^{2}$, $\omega_z = \omega_{\mathrm{cres}}$
    & 0.325$v$ & 0.1 & 0.220 & 109.7 & 77.1 & 5.83 & 228.3 & 305.4 & 11.12 \\[2pt]
  Coupled sinusoidal forcing\newline of $A_f = 23.8v^{2}$, $\omega_z = \omega_{\mathrm{cres}}$
    & 0.325$v$ & 0.1 & 0.200 & 110.0 & 77.3 & 5.93 & 234.4 & 311.7 & 9.29 \\
\end{longtable}
\endgroup

Since the drag increase due to angle of attack is proportional both to the square of $\alpha$ and
to the square of the airspeed, for instance, it follows that a transient disturbance producing a
given maximum angle of attack will reduce a model's overall altitude capability more when it
occurs near burnout (when the velocity is high) than near liftoff (when the velocity is low).
Because the corrective moment coefficient $C_1$ increases as the square of the airspeed, however, a
much stronger disturbance is required to produce a given maximum angle of attack near burnout than
near liftoff---and a disturbance of a given \emph{strength} occurring near burnout will produce a
reduction in performance roughly equivalent to that resulting from a disturbance of the same
strength occurring near liftoff. The quadratic dependence of drag upon angle of attack also means
that a disturbance of a given type occurring at a given time after liftoff will produce a
performance degradation roughly proportional to the square of its strength.

Given a set of in-flight disturbances of different types, all of which arise at the same time into
the flight and produce identical values of the maximum angle of attack in response, the severity
of the resulting performance degradations may be ranked according to the type of forcing function
involved as follows:
\begin{enumerate}[label=(\alph*)]
  \item impulse (least severe)
  \item step of constant intensity
  \item sinusoid
  \item step of intensity proportional to $v^{2}$ (most severe)
\end{enumerate}

A rocket that is spinning about its longitudinal axis will have a greater drag at zero angle of
attack than will the same vehicle if constructed so as not to roll. That the drag increase
associated with fitting a model with angled ``spin fins'' can be considerable has already been
shown by the analysis of Section~\ref{ch3:sec:5.3} of Chapter~\ref{ch3}. If, however (as has been
noted by Thomas McKim), the product of $I_R$ and $\omega_z$ is great enough to cause the
\emph{angular momentum} about the longitudinal axis to be very large, the maximum \emph{angle of
attack} due to a transient disturbance of a given strength will be significantly \emph{reduced}
over what it would have been without roll. During much of the response, therefore, the overall
drag of the rolling rocket may be less than the drag associated with the same vehicle built so as
not to spin. The response will, however, also \emph{persist} for a \emph{longer time} than it would
in the case of zero roll rate because the inverse time constant of the slower mode of oscillation
is reduced as the product $I_R\omega_z$ is increased. Both effects are due to the fact that the
presence of angular momentum about the roll axis \emph{causes the model to act as if it had a
larger longitudinal moment of inertia than its actual value of $I_L$} and thereby decreases the
amplitude, frequency, and decay rate of the response. Model rockets which are well designed to
begin with almost always suffer, rather than benefit, from the use of spin since the reduction in
maximum angle of attack due to any given disturbance is more than offset by the increased duration
of the transient response and the increase in drag at zero angle of attack associated with
spinning the model by aerodynamic means. There are a few exceptions to this rule; for example, the
use of spin fins and/or spin motors in the first stage of a heavy, multistaged model that might
otherwise suffer excessively from step or impulse forcing during its relatively slow liftoff. In
most cases where the incorporation of spin into a design improves its performance, however, it
will be found that it is because the design has not been optimized in the non-spinning
configuration.

Finally, the ability of a given model to resist a disturbance of a given type can be evaluated
from the following brief summary of the results of Chapter~\ref{ch2}:
\begin{enumerate}[label=(\alph*)]
  \item Resistance to step forcing increases with an increasing corrective moment coefficient
    $C_1$;
  \item Resistance to impulsive forcing increases with an increase in the product $I_L\omega$,
    where $\omega = \sqrt{\frac{C_1}{I_L} - \frac{C_2^{2}}{4I_L^{2}}}$; and
  \item Resistance to sinusoidal forcing increases with an increase in the product
    $C_1\zeta\sqrt{1 - \zeta^{2}}$, where $\zeta = \frac{C_2}{2\sqrt{C_1I_L}}$.
\end{enumerate}
The effect of a nonzero roll rate is to increase the apparent value of the longitudinal moment of
inertia, thereby lessening the angular frequency of the slower pitch and yaw oscillation modes and
decreasing the effective damping.
